\section{Introduction}

Ghana's forest cover faces accelerating degradation. Between 2001 and 2024, the nation lost approximately 1.2 million hectares of tree cover—a decline that threatens biodiversity, destabilizes local climate patterns, and releases substantial carbon emissions into the atmosphere. The drivers? Agriculture expansion dominates, accounting for roughly 90\% of all deforestation; cocoa and oil palm plantations carve through primary forests while shifting cultivation nibbles at the margins \citep{curtis2018classifying}.

This pattern is not unique to Ghana. Across West Africa, the tension between economic development and ecological preservation plays out in forest margins, where smallholder farmers and industrial agriculture compete for land against conservation imperatives. Yet Ghana presents a particularly stark case. The country hosts biodiversity hotspots—remnant patches of Upper Guinean forests—while simultaneously pursuing aggressive agricultural modernization. Policy oscillates between REDD+ commitments and agricultural intensification targets, rarely reconciling the two.

What complicates this further: forests themselves grow. Natural regeneration occurs; reforestation programs plant saplings. But these gains pale against losses. The data reveals an asymmetry. While Ghana added approximately 230,000 hectares of tree cover between 2000 and 2020, losses during the same period exceeded this five-fold. The net change trajectory tilts sharply downward.

Can we quantify the trade-offs? Economic models of forest management typically optimize a single objective—maximize timber revenue, minimize carbon emissions, or sustain agricultural output. Real-world decisions operate under multiple, conflicting pressures: farmers need income today; policymakers face carbon commitments tomorrow; conservationists advocate for biodiversity now. Dynamic optimization frameworks offer a method to navigate this complexity, balancing immediate agricultural returns against long-term ecosystem services under explicit constraints.

This paper constructs and solves such a model for Ghana. We formulate deforestation and reforestation as decision variables within a dynamic optimization problem spanning 100 years, where forest stock evolves according to logistic growth tempered by human intervention. The objective function incorporates agricultural revenue from cleared land, carbon pricing mechanisms applied to emissions and sequestration, amenity values from standing forests, and the direct costs of reforestation programs. Constraints ensure biological feasibility: you cannot deforest more than exists; reforestation cannot exceed ecological carrying capacity.

Our analysis proceeds in three stages. First, we document historical deforestation patterns using Global Forest Watch data, disaggregating losses by driver category and regional distribution. This empirical foundation establishes baseline conditions and validates parameter choices. Second, we specify the optimization model—a finite-horizon problem solved via Sequential Least Squares Programming—and derive optimal forest management trajectories under current economic parameters. Third, we conduct sensitivity analysis on carbon pricing, exploring how policy interventions might alter optimal decisions.

The findings challenge conventional wisdom. Under current carbon prices (\$20 per Mg CO$_2$e), the optimal strategy permits continued deforestation, albeit at reduced rates compared to historical trends. Only when carbon prices exceed \$45 per Mg does the model prescribe immediate cessation of deforestation paired with aggressive reforestation. This threshold sensitivity has direct policy implications: weak carbon markets cannot induce forest conservation through pricing alone; complementary regulatory mechanisms remain necessary.

We contribute to the literature in two ways. Methodologically, we integrate driver-specific deforestation data into an optimization framework, allowing policy analysis that accounts for heterogeneous land-use motivations. Empirically, we provide the first comprehensive assessment of Ghana's forest dynamics post-2020, extending previous analyses that terminated in 2015. The timing matters: cocoa prices surged between 2020 and 2024, potentially accelerating agricultural expansion into forest margins.

The paper unfolds as follows. Section 2 reviews the literature on forest management optimization and deforestation drivers in West Africa. Section 3 describes our data sources and details the dynamic optimization model. Section 4 presents results—both descriptive patterns from the data and normative prescriptions from the model. Section 5 discusses policy implications, particularly regarding carbon market design and agricultural intensification strategies. Section 6 concludes.
